\documentclass[11pt]{article}

\usepackage[margin=1in]{geometry}
\usepackage{amsmath,amssymb}
\usepackage{graphicx}
\usepackage{hyperref}
\usepackage{booktabs}
\usepackage{xcolor}
\usepackage{float}
\usepackage{microtype}
\usepackage{parskip}

\hypersetup{
  colorlinks=true,
  linkcolor=blue,
  urlcolor=black,
  citecolor=black,
}

\title{Ampleforth Network Durability}
\author{Fragments Research}
\date{October 2021}

\begin{document}
\maketitle

\begin{abstract}
Since its launch in 2009, Bitcoin has proved to be a resilient, censorship-resistant asset.
Like gold and many other physical commodities, Bitcoin does not rely on underlying collateral
or any central authority to maintain its state of distribution. However, also like gold, Bitcoin
experiences unit-price volatility and deflation, which undermines its use as a unit of account.
One solution is to separate an asset's float volatility, which can be fat-tailed and unbounded,
from that asset's unit price volatility, which can be mean-reverting. Launched in 2019, the
Ampleforth protocol has shown that this is possible and sustainable through extreme market
conditions.

\bigskip
AMPL is a fully-algorithmic unit of account developed on the Ethereum blockchain. Its
protocol transfers the volatility of demand from price per token to the number of tokens in
user wallets by automatically adjusting supply in response to its volume-weighted-average
price exchange rate. Holders of AMPL experience the high float volatility that characterizes
all digital assets. Between July 2020 and October 2021 the market cap of AMPL expanded by
a maximum of 84,736\% from trough to peak and contracted by a maximum of 93\%. However,
contracts denominated in AMPL have remained relatively stable in value through both secular
increases and decreases in demand over the same period. The price of AMPL has oscillated
around \$1.00, and remains within 20\% of that target 70\% of the time.

\bigskip
AMPL has additional dependencies which make it less resilient than BTC. Specifically, the
Ampleforth protocol requires the existence of well-functioning exchange markets between AMPL
and other assets. While the existence of these markets in 2017 and 2018 was not assured, it is
clear today that these markets are here to stay. By separating its stock volatility from its unit
volatility, AMPL allows the denomination of long-term on-chain contracts and the creation of
fully-decentralized crypto-backed derivatives.
\end{abstract}

\section{Rebasing Policy}

The Ampleforth Protocol targets the CPI adjusted 2019 US dollar and automatically expands
or contracts the quantity of tokens in user wallets based on price. When the
\texttt{price\_exchange\_rate} of AMPL $>$ \texttt{1 2019 USD} the quantity of tokens in user
wallets proportionally increases. When the \texttt{price\_exchange\_rate} of AMPL is $<$
\texttt{1 2019 USD} the quantity of tokens in user wallets proportionally decreases.

\begin{itemize}
  \item Supply adjustments, termed ``rebases'', are applied by updating a global scalar
    coefficient of expansion every day at \texttt{2AM UTC}. This adjustment is applied
    universally to all addresses and doesn't require any transaction between peers.
  \item Rebases only adjust supply if the price of AMPL deviates from its target by $>$ the
    \texttt{deviation\_threshold}. This threshold is set by an adjustable hyper-parameter and
    is currently set at 5\%.
  \item Rebases are smoothed by a \texttt{reaction\_lag} parameter. If the price of AMPL is
    $X\%$ above the target, the policy adjusts supply by
    $\left(\frac{X\%}{\text{reaction\_lag}}\right)$. This parameter is currently set to 10.
  \item The Ampleforth protocol's supply changes are proportional and non-dilutive. If a user
    owns $Y\%$ of the network before a rebase, the user will always own $Y\%$ of the network
    unless the user buys or sells more AMPL.
\end{itemize}

\section{Separating Price and Stock Volatility}

Price and stock volatility are often conflated, and for good reason. Most assets have an
inelastic number of units in their float. As a result, changes in their unit value produces
exactly proportional changes in their float value. Looking at Facebook's stock as an example,
if the stock price of a Facebook share were to increase by 10\% its market cap would
correspondingly increase by 10\%. The same is generally true of currencies. Ampleforth's
elastic supply policy, however, deliberately counteracts this. Recall that when the
volume-weighted-average unit price of AMPL trades above the protocol's price target, the stock
of AMPL increases proportionally \emph{until} price is restored to its equilibrium band and
vice-versa. Below, we show that the Ampleforth rebasing policy produces a unit price
distribution that is structurally different from the network's stock distribution. Notably,
AMPL's price volatility distribution reverts to its mean target price, while AMPL's stock
volatility remains fat-tailed.

\subsection{Price Volatility Distribution}

Here, we present the distribution of time spent across AMPL price deviations from its target
from December 2019 to October 2021. Although there is no theoretical upper bound to price,
the pattern demonstrates a clear average around the Ampleforth protocol's price target.

\begin{figure}[H]
  \centering
  \includegraphics[width=0.85\textwidth]{figures/price_deviation_dist.pdf}
  \caption{Distribution of AMPL price deviations from target (December 2019--October 2021).}
  \label{fig:price-deviation}
\end{figure}

\subsection{Stock Volatility Distribution}

Here we present AMPL's supply series over the same period (December 2019 to October 2021).
It can be seen below that, like a typical floating price token, AMPL's supply series (blue)
shows no clear long-run average, while its price series (red) hovers around a target price.

\begin{figure}[H]
  \centering
  \includegraphics[width=\textwidth]{figures/price_supply_overlay.pdf}
  \caption{AMPL price (red) and supply (blue) over time. The supply series exhibits fat-tailed
    behavior typical of floating-price tokens, while price reverts to its target.}
  \label{fig:price-supply}
\end{figure}

Finally we present the distribution of AMPL's supply changes over time. The distribution
takes as its inputs, rolling relative monthly supply changes. Below, it can be seen that the
distribution is asymmetric and fat-tailed, as is common for equities and floating price
cryptocurrencies.

\begin{figure}[H]
  \centering
  \includegraphics[width=0.85\textwidth]{figures/log_supply_change_dist.pdf}
  \caption{Distribution of AMPL rolling monthly supply changes showing asymmetric,
    fat-tailed behavior.}
  \label{fig:supply-dist}
\end{figure}

\subsection{General Comments}

The Ampleforth protocol's supply adjustments serve as an outflow for extreme swings in demand,
allowing AMPL's unit price to remain stable in the long-run. For this reason, contracts
denominated in AMPL remain similarly stable, even in the face of extreme market conditions
where other approaches are likely to require discretionary intervention or fail outright.

\section{Alternative Approaches}

The inevitable regulation of centralized stablecoins---and the rapid growth of decentralized
finance platforms that rely on them---has underscored the need for a decentralized unit of
account. To that end, many attempts have been made to produce decentralized stablecoins.
However few efforts, outside of AMPL, have elected to isolate this goal of creating a durable
unit-of-account from broader attempts to store near-term value. The key to AMPL's durability
is that it returns to a stable price by transferring the volatility of demand to supply,
rather than attempting to eliminate volatility. Systems that have attempted to remove
volatility altogether have thus far materialized as solutions that are either highly censorable
or highly unreliable.

The vast majority of alternatives fall under: coupon-based, crypto-collateralized, and
fractionally-collateralized approaches. Although exhaustive commentary on alternative systems
is beyond the scope of this document, we cover these three major categories below.

\subsection{Coupon Based ``Bond Token'' Approaches}

``Coupon coins'' peg their exchange rates by minting coins when their price is above the peg
and issuing interest-bearing coupons. Although these coupons are often referred to in project
whitepapers as ``bonds,'' they more closely resemble binary options that are paid off---if and
when a network needs to expand beyond its former all-time-high market cap in the future---or
not at all. Examples of this approach include Basis, BasisCash, Elastic Set Dollar, and
FEI---all of which have suffered major collapses.

Although these approaches, like AMPL, are fully algorithmic and do not rely on centralized
collateral; unlike AMPL they cannot survive secular decreases in demand. The key issue here
is that these systems cannot force a permanent reduction in supply. Once demand growth
flatlines after a correction the protocols create and magnify instability.

\subsection{``Crypto-Collateralized'' Approaches}

Although more promising than coupon-coins in theory, crypto-collateralized systems are
vulnerable to the volatility of underlying collateral. Most famously, DAI suffered a massive
collapse on Black Thursday when cryptocurrencies like ETH suffered rapid price declines. In
an emergency governance response, the protocol added USDC as collateral which now represents
$\sim$60\% of their collateral pool.

A key problem here is the use and reliance on fixed collateralization ratios. When such
collateralization ratios are held constant, the risk of liquidation is pushed off to market
participants in a manner that cannot be priced. So long as the network is over-collateralized,
stablecoins like DAI can maintain their price, but as a result it becomes the role of outside
markets to take on the black-swan risk of a price collapse in the underlying collateral assets,
and this type of risk is impossible to price. Many of these problems go away if the underlying
collateral asset is stable, as is the case with USDC, but introducing traditional assets
trades market risks for custodial and regulatory risks.

An alternative method of risk stratification has been developed by prl.one that transparently
re-segments the risk of underlying cryptocurrencies into tranches on-chain. Such derivatives
can eventually be used to produce on-chain safe-assets that serve as an alternative collateral
to USDC; and could therefore be used to salvage crypto-collateralized systems. But efforts on
this front are still in the early stages of development.

\subsection{``Fractionally-Collateralized'' Approaches}

Fractionally collateralized systems promise less than 100\% collateralization, noting that in
traditional finance we typically use capital more efficiently. All else being equal, this
provides less buffer for collateral collapse than over-collateralized systems, placing more
reliance on centralized collateral like USDC. Innovations in on-chain safe-asset collateral
are necessary for advancing these approaches.

\section{Applications Now and Looking Forward}

The decentralized finance movement aims to create an alternative financial ecosystem that is
open-source, borderless, and resistant to political tampering. AMPL enables the denomination
of stable on-chain contracts without any reliance on centralized custodians or buyers of last
resort. This capability propagates into a number of use cases including on-chain lending,
on-chain borrowing, and the creation of on-chain derivatives.

\subsection{Decentralized Debt Denomination}

The simplest AMPL use case is decentralized debt denomination. Typically when a person
borrows money they intend to put that money to work immediately and then pay it back later
with interest. However, when a loan contract is denominated with a floating price currency
like ETH, the borrower has to factor in the price volatility of ETH. To illustrate we can
follow a simple example:

\paragraph{ETH Denominated Loan Example.}
Imagine Alice borrows \texttt{1 ETH}, valued at \texttt{\$1000/ETH}, to be paid back at a
later date. If Alice sells the ETH to put \$1000 to work and the price of ETH increases to
\texttt{\$3000/ETH}, she will need to spend \$3000 to acquire the ETH in order to pay back
her loan.

Because loans denominated in price-volatile currencies are so risky, the majority of
borrowing activity on decentralized lending platforms is unfortunately denominated with
centralized stablecoins. AMPL eliminates this dependency and enables sustainable lending and
borrowing on decentralized lending platforms.

\paragraph{AMPL Denominated Loan Example.}
Imagine Bob borrows \texttt{1000 AMPL} to be paid back at a later date. If Bob sells the AMPL
to put \$1000 to work, he can be reasonably certain that the cost of re-purchasing the AMPL to
pay back his loan will be reasonably stable in the long-run due to AMPL's mean-reverting
price, even if the AMPL network profoundly increases or decreases in size.

Although AMPL holders experience the same kind of fat-tailed stock volatility expected of
floating-price tokens like ETH, contracts denominated in AMPL remain long-run stable, as is
the case with Bob's loan example above. AMPL is currently available as a lend and borrow
asset on AAVE.

\subsection{Decentralized Derivatives}

Derivatives are a special type of financial contract whose value depends on an underlying
asset. Such contracts require a unit of account for denomination and AMPL similarly enables
the creation of on-chain derivatives. To help illustrate this, let's explore a simple example:

\paragraph{AMPL Derivative Example.}
Imagine Bob wants to split the stock-volatility of AMPL into two derivative tranches, a
senior (low-risk) tranche and a junior (high-risk) tranche. Bob can author a simple smart
contract that accepts a redeemable AMPL deposit, and attributes $\frac{1}{10}$ of all future
supply changes to the senior tranche token while attributing $\frac{9}{10}$ of all future
supply changes to the junior tranche token. At the time of maturity senior tranches are first
in line to redeem followed by junior tranches second.

Bob will have created a senior derivative token that is considerably less volatile than the
underlying AMPL and a junior derivative token that is considerably more volatile, without any
added oracle risk.

\subsection{Decentralized ``Safe-Asset'' Collateral}

Following from the concept of on-chain derivatives above, correctly configured safe-asset
(senior) tranches can be used as transparent and robust collateral by the issuer of a
crypto-collateralized stablecoin. Because the tranche ratios of the respective derivatives
are pre-defined and redeemable, the volatility tolerance of a senior tranche token can be
clearly defined (i.e.: it can be initialized such that the safe-asset derivative can tolerate
an $X\%$ fall in demand). This sort of redeemable on-chain derivative does not require
liquidation markets and introduces no additional oracle risk. A breakthrough in safe-asset
collateral would allow crypto-collateralized stablecoins to reduce their reliance on tokens
backed by traditional assets (like USDC) as collateral, making them more decentralized.

\section{Oracle Risks}

To execute its automatic supply policy, the Ampleforth protocol accepts a
volume-weighted-average price exchange-rate through a network of Oracles. Although attacks on
this Oracle network can temporarily corrupt the inputs to Ampleforth's supply policy, in
theory, great care has been taken to negate the impact of Oracle Attacks. Most importantly,
because AMPL is non-custodial and the policy is proportional, no funds can be stolen or
redistributed by an Oracle Attack. To help illustrate this, let's walk through a simple
example.

\subsection{Stable Transfer Application Example}

Let's imagine a simple application that allows its users to send a dollar denominated amount
of money from one wallet to another, we'll call it Stable Transfer.

When Alice wants to transfer \$100 to Bob, the Stable Transfer contract queries an oracle for
the price-exchange rate of ETH to US dollars, and then it transfers an automatically
calculated quantity of ETH to Bob.

Since each transfer queries an oracle to calculate an ETH amount before transferring to its
receivers, oracle requests are $O(n)$ with the number of transfers. Each of these oracle
requests presents an exploit opportunity and incurs an added cost. Finally, if an attacker
exploits the exchange-rate oracle, or if there's a bug in the oracle, senders can transfer
the wrong amount of money resulting in a loss of funds.

\subsection{AMPL Example}

Alternatively, the AMPL equivalent of the Stable Transfer application, isn't really an
application at all. Alice would simply transfer 100 AMPL to Bob and there would be no oracle
query at the time of transfer. In the case of AMPL, rebases occur once daily and oracle
requests are $O(1)$ per day. More importantly, because the protocol's supply policy is
proportional and non-dilutive, oracle attacks cannot result in the theft or redistribution of
assets.

\section{Conclusion}

As mentioned above, floating price tokens like Bitcoin experience unit-price volatility and
deflation, which undermines their use as a unit of account. The Ampleforth solution transfers
the volatility of demand from price to supply using a simple, transparent, and non-custodial
policy. We've shown here that this policy separates AMPL's price volatility from its stock
volatility, allowing AMPL's price to revert to its target through both secular increases and
decreases in demand, without requiring centralized custodians, collateral, or buyers of last
resort.

\appendix
\section*{Appendix: Extended Analysis (March 2026)}
\addcontentsline{toc}{section}{Appendix: Extended Analysis (March 2026)}

Since the original publication of this report in October 2021, the Ampleforth protocol has
continued to operate uninterrupted through an additional four and a half years of market
activity. This period encompassed the collapse of multiple major centralized and algorithmic
stablecoin systems, a prolonged bear market in digital assets, and a subsequent recovery. The
extended dataset now spans December 2019 to March 2026, comprising 2,282 daily observations
compared to the 687 available at the time of the original analysis.

Over this extended period, the core finding of this report holds: AMPL's rebasing policy
continues to produce a mean-reverting unit price distribution that is structurally distinct
from its fat-tailed stock distribution. In fact, the price stability metrics have improved
with additional data. The price of AMPL now remains within 20\% of its target 83\% of the
time, up from 69\% over the original observation window. Similarly, the proportion of days
spent within 5\% of the target has increased from 25\% to 34\%, and within 10\% from 41\% to
57\%.

\begin{figure}[H]
  \centering
  \includegraphics[width=0.85\textwidth]{figures/appendix_price_deviation_dist.pdf}
  \caption{Distribution of AMPL price deviations from target (December 2019--March 2026).
    The distribution remains clearly centered around the target, with a tighter concentration
    than the original observation period.}
  \label{fig:appendix-price-deviation}
\end{figure}

The improvement in price stability is attributable to two governance-approved upgrades to the
protocol's rebase function. In 2022, the original linear rebase curve was replaced with a
sigmoid (S-shaped) curve, introducing asymptotic bounds that cap maximum supply adjustments
and limit exposure to oracle errors \cite{sigmoid-upgrade}. In 2024, the sigmoid curve was
further steepened and the deviation threshold was tightened from 5\% to 2.5\%, producing
faster mean-reversion by allowing the protocol to begin correcting at smaller price
deviations \cite{steepen-curve}. Together, these
upgrades have reduced the magnitude and duration of price deviations from the target,
producing a narrower distribution and a lower standard deviation (0.23 vs 0.35 over the
original period).

\begin{figure}[H]
  \centering
  \includegraphics[width=\textwidth]{figures/appendix_price_supply_overlay.pdf}
  \caption{AMPL price (red) and supply (blue) from December 2019 to March 2026. The supply
    series continues to exhibit fat-tailed behavior while price reverts to its target through
    multiple market cycles.}
  \label{fig:appendix-price-supply}
\end{figure}

The supply series continues to display the fat-tailed, asymmetric characteristics observed in
the original analysis. Over the full period, the market cap of AMPL expanded by a maximum of
30,668\% from trough to peak and contracted by a maximum of 99.7\%. Despite these extreme
movements in network value, the rebasing policy continued to absorb demand shocks and restore
the unit price to its equilibrium band without any discretionary intervention.

\begin{figure}[H]
  \centering
  \includegraphics[width=0.85\textwidth]{figures/appendix_log_supply_change_dist.pdf}
  \caption{Distribution of rolling monthly log supply changes (December 2019--March 2026).
    The distribution remains asymmetric and fat-tailed.}
  \label{fig:appendix-supply-dist}
\end{figure}

During this extended observation period, the protocol's supply expanded on 29\% of days,
contracted on 40\% of days, and remained unchanged on 31\% of days. The high proportion of
days with no supply change reflects periods where the price remained within the 5\% deviation
threshold, requiring no corrective action from the protocol.

These results reinforce the central thesis of this report: that a simple, transparent, and
non-custodial supply policy can produce a durable unit of account by transferring the
volatility of demand from price to supply. The Ampleforth protocol has now demonstrated this
property over more than six years of continuous operation, through multiple complete market
cycles.

\begin{thebibliography}{9}

\bibitem{sigmoid-upgrade}
Ampleforth Governance Forum,
``Rebase Curve Upgrade Proposal,'' 2022.
\url{https://forum.ampleforth.org/t/rebase-curve-upgrade-proposal/327}

\bibitem{steepen-curve}
Ampleforth Governance Forum,
``Steepen Rebase Curve to Tighten Price Distribution,'' 2024.
\url{https://forum.ampleforth.org/t/steepen-rebase-curve-to-tighten-price-distribution/772}

\end{thebibliography}

\end{document}
